\chapter{\IfLanguageName{dutch}{Realisatie en Implementatie}{Realization and Implementation}}%
\label{ch:realisatie}

In dit hoofdstuk wordt de technische en functionele realisatie van het proof of concept besproken. Waar de voorgaande secties zich richtten op de vereisten en het conceptuele ontwerp, focust dit deel op de implementatie aan de hand van de vooropgestelde requirements en wireframes. 
\\

De architectuur van de oplossing is fundamenteel opgebouwd uit drie samenwerkende lagen. Enerzijds is er de Salesforce In-App Guidance, die verantwoordelijk is voor de actieve begeleiding van de gebruiker over verschillende schermen en stappen heen. Anderzijds is er een setup assistant applicatie ontwikkeld: het Onboarding Dashboard. Dit dashboard, gebouwd als een moderne Lightning Web Component (LWC), fungeert als het controlecentrum. Het verdeelt de onboarding over kleinere stappen, stuurt de gebruiker naar de juiste onderdelen, en valideert op de achtergrond (via Apex) of de nodige data en configuraties succesvol zijn aangemaakt. Als laatste is er dan een documentatie-viewer, waarbij de gebruiker overal gelinkt is met de juiste documentatie voor de component waarin hij aan het werken is.
\\

Een belangrijk concept binnen deze setup assistant is de \textit{Auto-configuratie} functionaliteit. Om de tijd tot een werkende omgeving te verminderen, biedt het dashboard de mogelijkheid om met één druk op de knop al een aantal complexe configuraties en zelfs volledige Lightning Pages op de achtergrond te genereren. Hierdoor kan een gebruiker heel snel van start gaan met een echte omgeving en dan via de In-App Guidance verder gestuurd worden bij het zich eigen maken van het 24flow-platform.
\\

In de volgende paragrafen worden de belangrijkste pijlers van deze implementatie toegelicht: de dynamische opbouw van het dashboard, de integratie met de Salesforce guidance tools en de logica achter de documentatie-viewer.

\section{Architectuur van de front-end}

De visuele kern van het proof of concept wordt dus gevormd door het 'Onboarding Dashboard'. Omdat de nodige configuraties verspreid zijn over allerlei schermen, menu's en Salesforce-instellingen, is er gekozen om dit dashboard als centrum voor de onboarding te laten dienen. De gebruiker vertrekt steeds vanuit deze pagina om alles op te zetten en wordt ook steeds hiernaartoe teruggeleid.

\subsection{Opbouw van de User Interface (Dashboard)}
De interface is volledig gebouwd met Lightning Web Components (LWC) en het Salesforce Lightning Design System (SLDS). Het benutten van deze standaard architectuur is om twee redenen cruciaal voor de effectiviteit van het dashboard:

\begin{enumerate}
	\item \textbf{Consistente User Experience:} Door SLDS toe te passen, is de component visueel niet te onderscheiden van de standaard Salesforce-omgeving. Dit vergroot het vertrouwen en de herkenbaarheid voor de gebruiker.
	\item \textbf{Reactieve UI:} Wanneer de backend (Apex) registreert dat een configuratiestap is voltooid, of wanneer een proces loopt dat even moet laden (zichtbaar via de \texttt{lightning-spinner} component), worden de interface-elementen direct en asynchroon geüpdatet, zonder dat de gebruiker de pagina hoeft te herladen.
\end{enumerate}

Initieel in de wireframes was het dashboard zo opgebouwd dat aan de linkerkant de configuratiestappen stonden en rechts enkele extra tools.
Om de gebruiker echter zoveel mogelijk op het centrale scherm te houden en de documentatie hier ook rechtstreeks te kunnen tonen, is het dashboard in de uiteindelijke implementatie visueel opgedeeld in twee secties: een linkerkolom met de echte configuratiestappen en enkele tools daarboven en een rechterkolom puur om documentatie te tonen die steeds zoveel mogelijk met elke stap gelinkt is.

\begin{figure}[hbt!]
	\centering
	\includegraphics[width=0.9\textwidth]{DashboardOverview.png}
	\caption{Overzicht van het Onboarding Dashboard}
	\label{fig:dashboard}
\end{figure}

\subsection{De Setup-kolom (Configuratie \& Tools)}
Aan de linkerkant bevindt zich de interactieve onboarding component (\textit{24flow Setup Assistant}). Deze is opgebouwd uit de volgende logische lagen:

\begin{itemize}
	\item \textbf{Voortgang, Data \& Default configuratie:} Bovenaan bevindt zich een overkoepelende voortgangsbalk, waarop de gebruiker in percentages ziet hoeveel hij al geconfigureerd heeft. Deze balk is reactief: de logica combineert actuele database-controles uit de backend met lokale browser-opslag (\textit{localStorage}) voor handmatige taken, om zo het percentage te berekenen. Direct daaronder worden twee blokken getoond met algemene configuratie: 
	\begin{itemize}
		\item \textit{Demo Data Management:} Enerzijds een knop om via een Apex-proces automatisch demodata te genereren, zodat de gebruikersomgeving gevuld wordt met realistische productiedata. Anderzijds een placeholder knop om een eigen CSV-bestand in te laden met echte data van de klant, dit is een placeholder omdat het omvormen van deze data naar de Salesforce-omgeving buiten de scope van deze proof of concept valt. Zolang er geen data in de omgeving beschikbaar is worden de auto-generatie knoppen geblokkeerd. Onderaan de component is eveneens een reset-knop voorzien om deze data veilig te verwijderen en de gebruikersomgeving op te kuisen.
		\item \textit{Master Auto-Configuration:} Een geavanceerde functionaliteit waarmee de gebruiker voor alle 24flow-componenten een aantal default instellingen aanmaakt waar mogelijk en deze direct op automatisch gegenereerde pagina's plaatst. Hiermee kan een gebruiker heel snel een basisconfiguratie van 24flow in werking zien en testen.
	\end{itemize}
	
	\item \textbf{Factory Setup Complete (Conditioneel):} Een belangrijk moment in de onboarding is de fase direct na het uitvoeren van de bovengenoemde Auto-Configuratie. Om te voorkomen dat de gebruiker het overzicht verliest nadat de backend op de achtergrond allerlei instellingen en pagina's heeft gegenereerd, past de interface zich dynamisch aan. Er verschijnt een nieuw scherm ("Default Factory Setup Complete!") dat de gebruiker direct twee logische vervolgstappen aanbiedt:
	\begin{itemize}
		\item \textit{See how it works (Quick Tour):} Een knop die de gebruiker onmiddellijk door de zojuist gegenereerde schermen en instellingen leidt aan de hand van een overkoepelende walkthrough. Dit laat in realtime zien hoe de geconfigureerde onderdelen met elkaar samenwerken.
		\item \textit{Make it yours (Configure Manually):} Dit vestigt de aandacht  op de onderstaande hiërarchische takenlijst, waar de gebruiker de applicatie verder naar wens kan configureren en de verschillende standaardinstellingen kan bijwerken.
	\end{itemize}
\end{itemize}

\begin{figure}[hbt!]
	\centering
	\includegraphics[width=0.9\textwidth]{DashboardAutoConfig.png}
	\caption{Overzicht van het Onboarding Dashboard na auto-configuratie}
	\label{fig:dashboard_autoconfig}
\end{figure}

\begin{itemize}
	\item \textbf{Hiërarchische Takenlijst:} Daaronder volgt het stappenplan in een uitklapbare structuur. Configuratietaken zijn gegroepeerd met duidelijke titels, beschrijvingen, een indicatie van de geschatte uitvoeringstijd en een aanduiding of deze stap wel of niet gedaan is. Hierop bevindt zich ook een knop om voor deze specifieke stap automatisch terug de default configuratie te laten maken.
	
	\item \textbf{Substappen:} Binnen elke hoofdtaak bevinden zich enkele specifieke substappen met een geschatte uitvoeringstijd. Op dit niveau kan de gebruiker ervoor kiezen om een stap handmatig te configureren (wat de navigatie triggert en de walkthrough opstart), direct gerelateerde documentatie te openen, of de taak handmatig als voltooid te markeren.
\end{itemize}

\begin{figure}[hbt!]
	\centering
	\includegraphics[width=0.9\textwidth]{DashboardSubstep.png}
	\caption{Uitgeklapte hiërarchische takenlijst binnen het dashboard.}
	\label{fig:dashboard_open}
\end{figure}

\subsection{Documentatie-viewer}
Aan de rechterkant van het scherm bevindt zich de documentatie-viewer. Technisch gezien is deze component opgebouwd rond een \texttt{iframe}, waarmee de externe webpagina's van de 24flow-documentatie direct binnen de Salesforce-interface worden getoond. Door deze direct naast de setup-kolom te plaatsen, kan de gebruiker de uitleg uit de 24flow-documentatie raadplegen en tegelijkertijd links in het scherm de configuratie-knoppen bedienen.
\\ 

Technisch communiceren de twee kolommen via de \textit{Lightning Message Service} (LMS); zodra er links een stap wordt geopend, of de documentatieknop wordt aangeklikt, ontvangt de viewer een bericht en laadt direct de juiste URL van de respectievelijke stap in het \texttt{iframe}. Hiermee kan de gebruiker per configuratie stap direct de gelinkte documentatie openen in hetzelfde scherm. Dit elimineert de noodzaak om constant tussen verschillende browsertabbladen te wisselen en zelf op zoek te gaan naar de locatie van bepaalde documentatie, wat het vinden van de nodige info aanzienlijk versnelt. Als de gebruiker de documentatie toch wil gebruiken in een ander tabblad kan dit ook door op de pop-out knop te drukken aan de bovenkant van de Documentatie-Viewer.
\\

Om de documentatie mee te kunnen nemen doorheen de rest van de applicatie en niet enkel op het initiële scherm te tonen is er in de Salesforce utility bar aan de onderkant van het scherm steeds een knop met een venster waarin de Documentatie-Viewer ook te zien is. Hierbij is de effectief getoonde documentatie rechtstreeks gelinkt met de pagina waarop de gebruiker zich bevindt, wanneer dit bijvoorbeeld een orderpool pagina is zal deze documentatie te zien zijn.

\begin{figure}[hbt!]
	\centering
	\includegraphics[width=0.9\textwidth]{DocumentationViewer.png}
	\caption{De documentatie-viewer geopend via de Salesforce Utility Bar.}
	\label{fig:doc_viewer}
\end{figure}

\subsection{Interactieve Begeleiding (Walkthroughs)}
De visuele ondersteuning wordt afgemaakt door de integratie van Salesforce In-App Guidance. Waar het dashboard en de viewer het controlecentrum vormen, vindt de effectieve begeleiding plaats op de verschillende applicatiepagina's zelf.
\\

Zodra de gebruiker een taak start vanuit het dashboard, wordt de gebruiker verplaatst naar de juiste pagina voor die bepaalde configuratiestap. Hierbij start de gekoppelde walkthrough en is de eerste prompt zichtbaar, dit is ofwel een targeted prompt of een floating prompt afhankelijk van de situatie. Aan de hand van deze gekoppelde prompts wordt de gebruiker dan stap voor stap gestuurd bij het opzetten van zijn instellingen. De laatste prompt bevat steeds een link terug naar het dashboard en eventueel is er ook een docked prompt aanwezig met een video of andere multimedia indien dit beschikbaar is.

\begin{figure}[hbt!]
	\centering
	\includegraphics[width=0.9\textwidth]{WalktroughStepImpl.png}
	\caption{Start van een walkthrough met een targeted prompt (stap 1).}
	\label{fig:walkthrough_1}
\end{figure}
\pagebreak
\begin{figure}[hbt!]
	\centering
	\includegraphics[width=0.9\textwidth]{WalktroughStepImpl2.png}
	\caption{Vervolgstap binnen de walkthrough met een floating prompt (stap 2).}
	\label{fig:walkthrough_2}
\end{figure}


\section{Architectuur van de Backend}
De backend van het proof of concept heeft drie duidelijke kerntaken: het aanleveren van de inhoud voor het dashboard, het uitvoeren van de verschillende configuratie acties in de achtergrond en het implementeren van de documentatie-viewer.

\subsection{Custom Metadata Type}
\label{subsec:dynamische_datastructuur}
Om het dashboard overzichtelijk te houden, zitten de teksten, stappen en documentatielinks niet vast in de front-end code. In plaats daarvan wordt de structuur van het onboardingproces bewaard als data in een \textit{Custom Metadata Type}.
\\

Bij het inladen roept het dashboard deze configuratie op. De applicatie leest deze data (geformatteerd als JSON) uit om de stappenlijst op het scherm visueel op te bouwen. Door de inhoud en de weergave op deze manier te scheiden, is het mogelijk om later eenvoudig stappen toe te voegen of aan te passen zonder dat de code van de interface gewijzigd hoeft te worden.

\subsection{De Backend Controller}
\label{subsec:backend_controller}
Terwijl de front-end zorgt voor de visuele weergave, regelt een centrale Apex-klasse (\texttt{OnboardingController}) de effectieve bedrijfslogica en de communicatie met de database. Deze backend heeft vier concrete hoofdtaken:

\begin{itemize}
	\item \textbf{Data Ophalen:} Het inladen van de dashboard-opbouw vanuit de metadata, zodat de stappenlijst direct correct op het scherm verschijnt.
	\item \textbf{Statuscontrole:} De backend controleert met gerichte database-queries of taken daadwerkelijk zijn uitgevoerd. Hierdoor toont de voortgangsbalk altijd de actuele, ware status.
	\item \textbf{Automatische Configuratie:} Het systeem genereert met één actie de default waarden voor instellingen en stuurt de creatie van de benodigde UI-schermen aan.
	\item \textbf{Demodata Beheren:} Het genereren van realistische testdata om een lege applicatie te vermijden, inclusief de logica om de omgeving (zowel data als gegenereerde schermen) weer veilig op te kuisen.
\end{itemize}

\subsection{De Salesforce Metadata API}
Om de ware kracht van de \textit{Automatische Configuratie} te begrijpen, is het cruciaal om het onderscheid te maken tussen standaard Salesforce-data en 'metadata'. Waar standaard Apex-code moeiteloos datarecords (zoals productieorders of teams) kan aanmaken, is het van nature niet in staat om de gebruikersinterface zelf aan te passen. Om daadwerkelijk schermen te kunnen opbouwen, maakt het proof of concept op de achtergrond gebruik van een integratie met de \textbf{Salesforce Metadata API}.

\subsubsection*{De Meerwaarde van Dynamische Generatie}
Normaal gesproken vereist het opzetten van een visuele component, zoals de 24flow Team Cockpit, aanzienlijk veel handmatig werk. Een beheerder moet via de Lightning App Builder een nieuwe pagina creëren, de juiste componenten op het scherm slepen, de achterliggende configuratie-ID's handmatig koppelen, en vervolgens navigatietabbladen aanmaken en toewijzen. Dit proces is niet alleen tijdrovend, maar ook zeer foutgevoelig voor nieuwe gebruikers.

Door de inzet van de Metadata API wordt dit gehele proces gereduceerd tot één druk op de knop. De applicatie neemt de rol van de beheerder over en voert de volgende stappen autonoom uit, wat de opzettijd voor een nieuwe klant drastisch verkort:

\begin{enumerate}
	\item \textbf{Dynamische Pagina-opbouw:} De backend ontwerpt de visuele pagina (\textit{FlexiPage}) volledig in het geheugen. Hierbij plaatst het de benodigde Lightning Web Components op de juiste plek en injecteert het direct de zojuist gegenereerde instellingen in de interface. Hierdoor is het scherm onmiddellijk op maat gemaakt voor de gegenereerde data.
	\item \textbf{Slimme Foutafhandeling:} Het systeem is robuust ontworpen. Als een gebruiker een configuratie per ongeluk opnieuw genereert, detecteert het script dat de pagina al bestaat. In plaats van te crashen, zal het systeem de bestaande pagina en configuratie veilig updaten en overschrijven.
	\item \textbf{Directe Toegankelijkheid:} Zodra de pagina is opgebouwd, creëert de code onmiddellijk een navigatietabblad (\textit{CustomTab}) en injecteert deze in het menu van de 24flow-applicatie. De gebruiker kan het resultaat dus direct bekijken en gebruiken, zonder ooit de Salesforce-backend te hoeven openen.
	\item \textbf{Automatische Opschoning:} Deze automatisering werkt in twee richtingen. Wanneer de gebruiker besluit de demodata te wissen, gebruikt het systeem de API om niet alleen de data, maar ook de gegenereerde schermen en tabbladen permanent te verwijderen. Zo wordt de klantomgeving altijd netjes en in originele staat achtergelaten.
\end{enumerate}